\section{Purpose, conclusion and evidence rules}
\label{hdc-purpose}

This dossier transfers a sustained attempt to complete a rigorous
existence argument for a nontrivial four-dimensional Yang--Mills
quantum theory. The project starts from the repository
\emph{Yang--Mills-LLM-experiment}, hosted at
\texttt{ethanwsussman/Yang--Mills-LLM-experiment} with a README
note signed LA, and from the Ba{\l}aban renormalization literature. The detailed
baseline assessment below distinguishes the repository's stated
ambition from the precise scope of its displayed results.

At handoff there is a substantial local ultraviolet construction
in the working proof chain, together with new exact statements
about full finite Wilson measures. There is not yet a completed
continuum existence proof. A complete-density multiscale theorem,
observable estimates, and continuum identification still require
mathematical arguments.

The final research step, Notebook 39, completes a restricted
normalized replacement operation on arbitrary collections of
separated fixed boxes in the actual bare Wilson theory. Its
denominator is controlled, its exterior marginal is preserved,
and its effect on local bounded observables does not grow with
the number of distant boxes. This is a natural stopping point
for the present round: one finite-measure operation is complete,
and its remaining distance from a multiscale construction is explicit.

\subsection{What a proof claim in this bundle means}

The documents distinguish five kinds of statement.
\begin{enumerate}
\item A \emph{derived working result} has a written argument in the
included notebooks, with its assumptions and scope stated. Many
recent modules have separate internal mathematical audits.
This is not independent professional verification of all earlier
dependencies.
\item A \emph{conditional reduction} proves an implication from
explicit model bounds. Its hypotheses are not counted as established
merely because the reduction is proved.
\item A \emph{published input} is tied to an actual primary statement.
Its use requires matching its domains, normalizations and uniformities.
\item An \emph{open interface} is a step of this proposed route
for which no complete argument is supplied here. A different
route to existence might replace it.
\item A \emph{superseded argument} is retained for provenance, with
its correction or replacement identified. Its presence in the
research archive is not an endorsement.
\end{enumerate}

The terms our addition and derived here identify contributions
to this project record. They are not claims of novelty relative
to the whole literature. The compact-group Laplace, tree-gauge,
coupling and chessboard mechanisms have standard antecedents.
Their precise assembly and use in the present program are what
this dossier records.

\subsection{The archive's order of authority}

Read the final handoff documents for the current map, then the
numbered notebooks for statements and proofs. Read each
notebook's dependencies before treating its theorem as
unconditional. The research archive contains fragments,
calculation notes, audits and historical checkpoints. Later
corrections override earlier speculation and earlier next-task
paragraphs.

The catalogue supplies an entry for every numbered notebook:
its location, central results, assumptions, dependencies and
limitations. The manifest identifies exact files by SHA-256.
The repository is pinned separately; a future remote update
is not silently substituted for the examined commit.

\section{The target and the hierarchy of conclusions}
\label{hdc-target}

The working regulator is a finite four-dimensional periodic
Wilson lattice gauge theory, principally for \(SU(2)\) in the
inherited construction. At positive coupling its compact group
integral exists automatically. The hard question is a nontrivial
continuum theory at vanishing lattice spacing, with the
quantum-field properties needed to identify it as Yang--Mills.

Three conclusions must be distinguished. A continuum limit of
separated scalar action-density correlations is a useful
observable sector. Reconstruction from a sufficient Euclidean
hierarchy could give a local relativistic quantum observable
theory. Identifying it as nontrivial Yang--Mills requires an
adequate gauge-invariant algebra and its relation to the lattice
theory and short-distance normalization. A physical mass gap is
a further spectral assertion. A result for \(SU(2)\) does not
by itself cover every compact simple gauge group.

Subsequential existence is a legitimate target. A unique cutoff
limit or preferred contact-term convention is not imposed
solely for convenience. What is indispensable is one consistent
limiting hierarchy with the requisite positivity, covariance,
regularity and physical interpretation.

\subsection{Distinct parameters and uniformities}

In the main local chain, \(M\) denotes preceding homogeneous
classical depth; \(\ell\) is the fixed primitive ratio;
\(R=\ell^s\) is a fixed composite quantum ratio; \(J\) is a
fixed auxiliary grouping; \(N\) is a finite period; \(b\) is
a small analytic radius; \(g\) is the running coupling; and
\(K\) is a large endpoint ratio in determinant comparisons.
In physical lattice statements \(\eta\) denotes mesh size.
Individual theorem statements control any changes of notation.

Uniformity in \(M,N\) on a fixed small ball is not uniformity on
arbitrary exterior fields. Constants depending on fixed \(R,J\)
are not uniform as either grows. A fixed-box bound uniform in
surrounding volume is not a bound on a physical box containing
more and more lattice sites. An infinite structural kernel is
not automatically the actual finite winding kernel. Required
old-depth shifts, endpoint ratios and source derivatives must
belong to one joint structural limit.

\section{The local construction added in this project}
\label{hdc-local}

The main route is
\[
 \begin{gathered}
 \text{flat coercivity and classical minimizers}\\
 \longrightarrow\text{actual local quantum integral}\\
 \longrightarrow\text{Ward extraction and finite induction}\\
 \longrightarrow\text{actual positive determinant drift}\\
 \longrightarrow\text{arbitrarily long finite weak local histories}.
 \end{gathered}
\]
The catalogue supplies every intermediate notebook, rather than
replacing this chain by an appeal to perturbation theory.

\subsection{Flat and nonlinear classical analysis}

Notebooks 06--17 develop flat coercivity, constrained lifts,
weighted inverse kernels, refinement limits, pointwise regularity
and nonlinear response estimates for the homogeneous classical
minimization problem. The construction tracks physical pairings,
contour constraints, gauge directions, and the relation between
a prescribed coarse field and its fine minimizer.

Their endpoint is a family of local minimized classical actions
\(S_M\), analytic charts, and localized response kernels near
the trivial orbit. The support estimates matter because quantum
integration must preserve spatial summability, not merely
analyticity on a finite-dimensional ball.

These results concern regular homogeneous backgrounds. They do
not give a minimizer theorem on every determining set with
holes, arbitrary exterior boundary fields, or independent
complex currents. Their estimates cannot be used on those
domains by renaming variables.

\subsection{The actual local quantum integral}

Notebooks 18--28 formulate and control one fixed composite
ordinary coarse map \(F_R=\mathcal M_\ell^s\), its chosen
direct-\(R\) contour gauge, and its retained-coordinate
small-field sector. The nonlinear constraint chart
\(\Phi_V(z)\) keeps the retained coordinates unchanged and
solves the pivot variables. Its integration field is
\[
                 U=e^{i\Phi_V(g\xi)}V,
\]
not \(e^{ig\Phi_V(\xi)}V\). This distinction affects
higher-order interactions and the exact Jacobian.

The pullback includes Haar densities, the square ambient
straightening, the pivot constraint derivative, the contour
gauge penalty, the actual reduced Hessian, and the independent
retained link/color-ball window. A rectangular determinant of
the retained embedding cannot replace the constraint Jacobian.

The compact gauge insertion uses the exact normalizer
\[
 c(g)=\int_G\exp[-(1-\operatorname{Re}\tau h)/g^2]\,dh.
\]
For a literal represented input density the local output
has the following factorization, with notation abbreviated
but the normalization displayed:
\begin{align*}
 &Z_{\rm in}\,c(g)^{-n_s}h_0^{n_e-n_c}D_R(V)\,
 g^{n_z}(2\pi)^{n_z/2}\det A_M(V)^{-1/2}\\
 &\quad{}\times e^{-S_M(V)/g^2+E_{\rm raw}(V)}
 \mathbb E_{A_M(V)^{-1}}
       [\chi_{\rm window}\chi_{\rm extra}e^{I(V,g,\xi)}].
\end{align*}
Here \(n_z=(\dim\mathfrak g)(n_e-n_c)\), \(D_R\) is the
actual pivot factor, \(h_0\) is the Haar density at the origin
in the chosen Lie coordinates, and \(I\) contains all higher
action, Haar, straightening and inherited fluctuation terms.
Individual notebooks define every factor without shorthand.

The common-domain argument prevents uncontrolled shrinkage of
the analytic field ball at every step. It uses the actual
pre-gauge representative and transports the \emph{entire}
inherited insertion, including its background value.
It does not merely change a fluctuation variable and forget
its cutoff. Strict connected support families and
overlapping-carrier expansions provide the spatial norms.

Composition of primitive constraints identifies the ordinary
coarse map. It does not identify one composite retained window
with an intersection of sequential primitive windows. Also a
holomorphic continuation of a local integral is not a
representation of the full positive measure outside its
real integration domain.

\subsection{Ward extraction and finite local induction}

Notebooks 29--31 use Ward identities of the total functional
to isolate its marginal Maxwell contribution. They do not
assume that each arbitrarily selected activity is separately
gauge invariant. The remaining interactions are relocalized
and controlled in a suitable representative of the exact
finite total.

That representative need not be the original canonical cluster
array. Equivalence of the total and its compensating vacuum
constant are tracked: changing \(E\) to \(E+c\) requires
changing \(Z\) to \(Ze^{-c}\). Dropping that factor would
destroy a later full-measure normalization comparison.

The finite induction bounds the local remainder for permitted
weak coupling sequences and sufficiently large compatible
periods. It does not itself force the couplings to remain
weak. Positive drift and finite shooting supply that part.

\subsection{The actual logarithmic drift}

Notebooks 32--37 compute the leading coefficient for the
actual constrained determinant. An exact soft reduction
expresses it as a fine scalar/vector pair and coarse factors.
The coarse Maxwell contributions and both heat-integral ends
are bounded uniformly in preceding depth.

On the remaining band \(1<s<K^2\), source-localized resolvents
remove soft penalties with an integrable error. Nonlinear
Ward jets through field degree four compare the actual
effective-action vector base with the Wilson base;
left-relative chart-current terms remain present.
A direct Wilson-to-continuum comparison controls the
remaining error and fixes signs and color multiplicities.

With \(\tau(H^2)=1\) and the full adjoint charge trace,
\[
 C_H=\sum_{q\in\operatorname{spec}(\operatorname{ad}_H)}q^2,
 \qquad \alpha_g=\frac{C_H}{48\pi^2},
\]
the working theorem is
\[
       \beta(\mathcal D_{M,K})=22\alpha_g\log K+O(1).
\]
The error is uniform in old depth and every applicable
joint infinite structural system, under the classical and
local quantum inputs explicitly used in the series.
For \(\tau=\operatorname{Tr}/N\), the color calculation
gives \(C_H=2N^2\); this does not establish earlier
construction inputs for every \(SU(N)\).

Here \(\beta=2[t^2]\) on the normalized affine probe.
The signed contributions of the mean-zero scalar and vector
factors to \(\mathcal D\) have coefficients
\(2\alpha_g\) and \(20\alpha_g\). An older point-pinned
ghost decomposition distributes logarithms differently;
its scalar coefficient cannot be combined with the new
mean-zero vector factor.

\subsection{Finite weak histories and the recursion actually tuned}

Finite moment transfer converts the positive margin in every
joint structural limit to a uniform finite block margin.
The geometric ratios are fixed first; a block of \(q\)
steps then dominates the bounded comparison error; periods
and the weak ceiling are chosen so the actual errors do
not consume that margin.

A continuous clipped finite recursion is shot to a fixed
weak endpoint \(\Gamma_*\), where \(\Gamma=g^{-2}\).
Positive future blocks show clipping did not occur on
the selected trajectory. With \(m\) blocks,
\[
 (\Gamma_*+mqD_{\rm abs})^{-1/2}
 \le g_{\rm bare}\le
 (\Gamma_*+m\delta/4)^{-1/2}.
\]
This supplies arbitrarily long finite \emph{homogeneous local}
weak histories. It does not claim uniqueness, consistency
between horizons, or that a full large-field density has
the same extracted recursion. That identification remains open.

\section{The full-measure additions}
\label{hdc-fullmeasure}

Notebooks 38--39 add exact finite positive-measure results.
Their bare Wilson probability theorems are independent of
the earlier determinant work. The local-germ embedding uses
the already established coarse chart and Ward identities.

\subsection{Exact positive history operators}

Fix a measurable equivariant completion of the actual local
coarse map, using its original mean on its good set and
its original reference string outside. This specifies a
global coarse observable of the original fine measure;
it does not invent a global analytic effective action.

The compact based-gauge insertion gives a probability
kernel preserving the coarse field. An actual retained
sheet yields a local density with
\[
                  F_*\mu\ge 1_D\ell_\mu\,m_C.
\]
For a nonnegative cutoff \(\theta\) supported in \(D\),
let \(m_\theta(W)=\int\theta(W^v)\,dv\). Define
\[
 \mathcal A_g\mu(dW)=
 \frac{\int\theta(W^v)\ell_\mu(W^v)\,dv}{m_\theta(W)}\,m_C(dW),
\]
with value zero when \(m_\theta=0\). The invariant
absolutely continuous density of the true pushforward
bounds every numerator orbit value. Hence
\[
       0\le\mathcal A_g\mu\le F_*\mu,\qquad
       \mathcal R_g=F_*-\mathcal A_g\ge0.
\]
These linear invariant operators retain singular parts
in the remainder. Their finite products account for
every history, raw mass and observable numerator exactly.

Local-germ equality uses a finite orbit-overlap lemma
and the literal invariant-ball cutoff Ward identity.
Its radius may shrink with graph diameter. Every
reconstructed real field must also lie where the input
measure actually has the claimed density.
Holomorphic continuation alone supplies neither condition.

\subsection{Actual Wilson rarity and conditional normalization}

For site-plane cells of fixed side \(b\ge3\) and torus
periods \(n_i=2m_i b\), exact Haar slab gluing and even
Schatten H\"older give a chessboard inequality.
A lower bound on the \emph{full} partition function yields
\begin{align*}
 \mu_g\left(\bigcap_{j\in S}B_j(\delta)\right)
 &\le z_b(g,\delta)^{|S|},\\
 z_b&=\min\{1,e^{-\delta^2/g^2+
       b^d(d\dim G\log(1/g)+C_{G,d})}\}.
\end{align*}
Each \(B_j\) forces an interior plaquette energy
\(1-\operatorname{Re}\tau U_p\ge\delta^2\).
Finite coloring gives prescribed-plaquette bounds with
exponent loss \(Q_b=b^d\binom d2\).
This is a normalized interacting Wilson estimate,
uniform in surrounding volume.

For a fixed box \([0,R]^4\) with small prescribed
boundary edges, a forest rooted at every boundary
vertex removes the interior gauge redundancy exactly.
At identity boundary its representative has a unique
zero-action field and a positive Hessian. Compactness
and the implicit-function theorem give a unique global
conditional minimum for nearby data. Numerator and
denominator have the same local factor \(g^n\);
the outer compact region has a fixed energy gap. Thus
\[
      \mu_{V,g}(W_{V,g,T}^c)\le Ce^{-cT^2},
      \qquad 1\le T\le r/g.
\]
The window uses an adjoint-invariant edge-max norm
about that conditional minimizer.

The boundary complex is \(S^3\). Boundary tree loops
have finite boundary-plaquette fillings, so small
boundary curvature supplies a boundary-only small
gauge. Conditioning on the exterior of vertex-disjoint
boxes makes their interiors independent. Combining
the estimates gives
\[
 \mu_g\left(\bigcap_{i=1}^m D_i\right)
 \le \left[Ce^{-cT^2}
       +24R^3 z_b(g,\varepsilon_R)^{1/Q_b}\right]^m.
\]
Here \(D_i\) is a bad boundary or a failed local
gauge/minimizer window. It is an event in the complete
bare measure. It is not yet the nonlinear averaged
constraint window of the iterated map.

\subsection{The final normalized replacement}

Notebook 39 conditions on the common exterior \(\omega\).
The interior laws factor as \(\prod_iQ_i^\omega\).
On a good boundary put
\[
 p_i(\omega)=Q_i^\omega(W_i^c)\le a_g,\qquad
 P_i^\omega=Q_i^\omega(\,\cdot\mid W_i).
\]
On a bad boundary take \(P_i^\omega=Q_i^\omega\) and
\(p_i=0\). The operation does not remove bad boundaries.

For the original exterior marginal \(\rho\), set
\[
          \nu(d\omega,dx)=\rho(d\omega)\prod_iP_i^\omega(dx_i).
\]
This preserves probability and, in the unnormalized
version, the original raw partition mass.
Its Radon--Nikodym factor is a product of good-window
indicators divided by their actual probabilities
\(1-p_i(\omega)\ge1-a_g\).

An explicit coupling changes an interior only on a
conditionally independent failure of probability \(p_i\).
For the exterior jointly with a selected set \(A\),
total variation, with convention
\(\sup_E|\mu(E)-\nu(E)|\), is exactly
\[
 \mathbb E_\rho\left[1-\prod_{i\in A}(1-p_i)\right]
 \le |A|a_g.
\]
The estimate is independent of the number of other
boxes. The full theorem gives sparse discrepancy
bounds, retains correlated bad boundaries, and
transfers bounded local sources with the same
normalized laws.

This completes a finite localized replacement theorem.
It does not cover arbitrary growing components,
repeated scales, or nonlinear coarse-average fibers.
Preserving normalization does not establish that the
homogeneous coupling recursion remains unchanged.

\section{The remaining theorem program}
\label{hdc-open}

\subsection{H2 Complete-density transfer for the chosen map}

The first major unresolved theorem is propagation of
a complete density for the same full Wilson measure
and coupling recursion as the local construction.
There are two routes.

An import route verifies the entire induction class
of Ba{\l}aban's complete-density theorem for the chosen
coarse map and operations. That theorem is conditional
on its own effective couplings staying weak.
Matching their expected size or one-loop sign is not
a match of the recursion. The primary-source map
appears in the baseline assessment and Notebook 38.

A direct route builds a spatial determining-set
decomposition with collars, a regular constrained
background, rough fluctuations and exceptional
regions. It proves exact good-component integration
with the actual window and Haar, gauge, pivot and
vacuum factors. It then bounds normalized replacement
on every allowed complementary component, with
constants summable in geometry and scale age.
Notebook 39 supplies a restricted fixed-box example.

The proof should state admissible exterior fields,
transition or holonomy data, constraints, conditional
measures, analytic domains and the component norm.
Its denominator lower bound must be compatible with
the numerator's boundary energy. An unrelated
Gaussian normalization cannot be inserted after
conditioning.

Several shortcuts are excluded. Unfixed raw link
marginals are Haar. Flat periodic fields can have
nontrivial holonomy. Even plaquettes of size
\(O(\eta^2)\) can have alternating physical curvature
whose adjacent norms prohibit every uniformly
bounded \(C^{1,\alpha}\) gauge. Small curvature
alone is not a smooth-background construction.

\subsection{H3 Observable and source estimates}

After complete-density control, external source
estimates are needed for a specified gauge-invariant
observable algebra. Analyticity in a connection
background is not analyticity in an external
composite-observable source. The latter changes the
density, local counterterms, translation properties
and possibly its large-field estimates.

A useful first target is renormalized action density
on reflected separated test supports. Its source
neighborhood must be uniform in mesh and volume
in the exact cells used by reflection. Mixed
derivatives must obey all-order distributional
bounds, for example factorial bounds in stated
test-function seminorms. Fixed-point-number bounds
with no controlled growth in that number do not
supply field reconstruction estimates.

Numerator and denominator must belong to the same
positive regulator construction. Notebook 39's
bounded-source theorem does not control renormalized
fields whose supremum norms grow with refinement.

\subsection{H4 A common continuum hierarchy}

Once actual model bounds are available, extract one
diagonal subsequence for a countable dense test
algebra and extend by uniform continuity. Products,
mixed insertions, subtraction conventions and
required contact extensions must be compatible
on that same subsequence. Compactness of analytic
interactions on a background ball is a different
problem from tightness or distributional
compactness of physical correlations.

Separated correlations may not determine every
distribution needed for local fields. A route
avoiding particular coincidence extensions should
state why; otherwise those extensions and
counterterms must be constructed compatibly with
positivity and covariance.

\subsection{H5 Positivity, rotations and reconstruction}

Reflection positivity must descend from the same
positive finite measures and limiting observable
algebra. Estimates for a substitute measure or
unrelated subsequence cannot simply be combined
with it.

Hypercubic symmetry is not \(O(4)\) covariance.
The reflection-based restoration mechanism has
conditional geometric and source-cell reductions
in the early notebooks. Its uniform model
estimates still need to hold for all transformed
supports, cell positions and all-order constants
used by that argument.

Osterwalder--Schrader reconstruction requires the
chosen hierarchy's precise regularity and
positivity hypotheses and an appropriate vacuum
sector. Unique vacuum or clustering needs its
own estimate. If existence can use a suitable
sector or decomposition instead, that choice
should be explicit.

\subsection{H6 Nontrivial Yang--Mills identification}

A positive beta coefficient does not establish
nonzero positive-energy states in the reconstructed
theory. A quantitative lower witness must survive
the same limit. In a reconstructed theory, a
strict inequality
\[
 \|\psi\|^2-\langle\psi,e^{-tH}\psi\rangle>0,\qquad t>0,
\]
would certify positive-energy content. The task is
to prove such a witness from the cutoff theory.

A non-Gaussian scalar composite is not by itself
an interaction witness: a quadratic composite of
a free Gaussian field is already non-Gaussian.
The algebra, tensor laws and short-distance
normalizations must be sufficient to identify
the limit as Yang--Mills. A sufficient
identification criterion should be stated and
proved; no speculative criterion is promoted to
a theorem in this dossier.

\subsection{H7 Physical mass gap and group scope}

The gap in an auxiliary constrained Gaussian,
typically at the blocking scale, is not a
positive gap in the interacting Hamiltonian
above the vacuum. A physical infrared estimate
must be uniform in the same continuum sequence.
An \(SU(2)\) construction does not settle the
all-compact-simple-group target. These are
separate from finishing the ultraviolet program.

\section{How a successor can use the bundle}
\label{hdc-continuation}

\subsection{A recommended first pass}

Read the start document, the baseline assessment
and catalogue entries 31--39. Then inspect
Notebook 37's actual coefficient theorem and
trace its upstream inputs, or independently
begin with Notebooks 38--39's bare Wilson
arguments. These routes need different
expertise and can be reviewed in parallel.

For the drift, prioritize relative
Hessian/current terms, nonlinear Ward jets,
the determinant decomposition, infrared and
source-localized traces, and the transfer from
joint structural limits to finite drift.
For complete-density transfer, prioritize
domains, coarse-map completion, retained-window
equality, exterior normalization and coupling
extraction.

\subsection{A concrete next research deliverable}

The next useful deliverable is a theorem for a
collared component with the actual averaged
constraint. It should specify:
\begin{enumerate}
\item the component, exterior data,
gauge/transition variables and compact
conditional integral;
\item the constrained background and retained
map with its full Jacobian;
\item an exact normalized replacement or
good/remainder identity preserving the fine
measure when components are summed;
\item a bound in the complete-density
induction's actual norm;
\item compatibility with neighboring regions,
scale changes, sources and coupling extraction.
\end{enumerate}

Verifying the published induction class may
be more efficient than rebuilding its
geometry. Choose by an exact hypothesis
comparison. A direct construction changing
the coarse map must redo its local
identification and drift transfer.

\subsection{Review and reproducibility}

Internal audits and successful compilation
are recorded, not called external
certification. A successor should look for
hidden domain enlargement, dependence on a
supposedly uniform parameter, an unrelated
conditional denominator, or identities
combined from different limits. These are
specific failure modes encountered and
repaired during this project.

Portable sources contain the assembled
notebooks and historical fragments.
Reproduction instructions give compile
commands; the manifest records delivered
bytes. Third-party articles are linked
rather than republished. The repository
can be fetched at its pinned commit and
checked against the recorded tree.

The outcome is a detailed foundation and
continuation program. The ambition remains
existence; the bundle makes the difference
between that ambition and the established
arguments inspectable.
